\notepage{54}{The Rydberg energy and the Bohr-radius correspondence}
The natural Coulomb energy scale is the Rydberg,
\[
 R_y=\frac{me^4}{2\hbar^2}\simeq13.6057\ \mathrm{eV},\qquad
 E_n=-\frac{R_y}{n^2}.
\]
The handwritten notes quote the numerical value
$13.605\,692\,53(30)\ \mathrm{eV}$ and attribute it to NIST; this is
retained as a historical quotation, not as a current precision value.

\smallheading{Radial asymptotics}
The source points to textbook pages 356--357 for explicit
eigenfunctions, without identifying the textbook. Their large-radius
behavior is
\[
 R_{n\ell}(r)\propto r^{n-1}\ee^{-r/(na_0)}\qquad(r\to\infty),
\]
whose power and exponential are independent of $\ell$ at fixed $n$.
Indeed,
\[
 \lambda=\frac{\sqrt{-2mE_n}}{\hbar}
 =\frac{1}{n(\hbar^2/me^2)}=\frac1{na_0},\qquad
 a_0=\frac{\hbar^2}{me^2}.
\]

\smallheading{Connection with the Bohr model}
For the largest orbital angular momentum at a given $n$, take
$\ell=n-1$. In that case the radial polynomial has no additional
nodes, and
\[
 \Psi_{n,n-1,m}(\mathbf r)\propto
 r^{n-1}\ee^{-r/(na_0)}Y_{n-1}^{m}(\theta,\phi).
\]
The probability of finding the electron in a spherical shell of radius
$r$ and thickness $\dd r$ is
\[
 \left(\int\abs{\Psi(\mathbf r)}^2r^2\dd\Omega\right)\dd r
 \propto r^{2n}\ee^{-2r/(na_0)}\dd r.
\]
Thus the radial probability density has its maximum at
\[
 \frac\dd{\dd r}\left(r^{2n}\ee^{-2r/(na_0)}\right)=0
 \quad\Longrightarrow\quad
 \boxed{r_{\max}=n^2a_0.}
\]
This equals the $n$th Bohr-orbit radius. The statement concerns the
maximum of a radial probability distribution, rather than a definite
electron trajectory.
